\editionPage[\pdfbookmark[0]{Démonstrations}{bm-preuves-debut-parent}]{preuves-debut}{Comprendre les preuves}{Pourquoi on peut ajouter des fractions}{Les arguments ci-dessous éclairent des propriétés utilisées dans le livre.}
\begin{blocEdition}{preuves-debut}
\editionEtape{Garder la même unité}Pour réunir des parts, elles doivent avoir la même taille. Sur une unité partagée en six, une moitié représente trois parts et un tiers deux parts.\label{fig-preuves-debut-cours-somme-fractions}\figureGeo{cours-somme-fractions}
\editionEtape{Compter}Trois sixièmes et deux sixièmes donnent cinq sixièmes. Le dénominateur indique toujours le partage ; seul le nombre de parts change.\editionEtape{Généraliser}Pour un même dénominateur $b>0$, $a$ parts et $c$ parts donnent $a+c$ parts : $\frac ab+\frac cb=\frac{a+c}b$.\editionRetenir{Pour des dénominateurs différents, changer d’écriture avant de réunir les parts. L’égalité des fractions assure que leur valeur n’a pas changé.}
\end{blocEdition}
\editionPage{preuves-2}{Comprendre les preuves}{La somme des angles du triangle}{Les arguments ci-dessous éclairent des propriétés utilisées dans le livre.}
\begin{blocEdition}{preuves-2}
\editionEtape{Une ligne auxiliaire}Tracer par C la parallèle à (AB).\label{fig-preuves-2-cours-triangle-angles}\figureGeo{cours-triangle-angles}
\editionEtape{Transporter les angles}Les angles formés en A et B se retrouvent en C grâce aux angles alternes internes égaux entre parallèles.\editionEtape{Former un angle plat}Ces deux angles et l’angle intérieur en C forment un angle plat. Leur somme vaut donc $\res{180^\circ}$.\editionRetenir{La ligne ajoutée transforme une question sur trois sommets en une question au même sommet. La preuve vaut pour tout triangle non aplati.}
\end{blocEdition}
\editionPage{preuves-3}{Comprendre les preuves}{Une médiane partage l’aire en deux}{Les arguments ci-dessous éclairent des propriétés utilisées dans le livre.}
\begin{blocEdition}{preuves-3}
\editionEtape{Identifier les bases}M est le milieu de [AB] : AM et MB ont la même longueur.\label{fig-preuves-3-cours-mediane}\figureGeo{cours-mediane}
\editionEtape{Identifier une hauteur commune}Les triangles ACM et BCM ont leurs bases sur la même droite (AB) et le même sommet C. Leur hauteur relative à cette droite est donc identique.\editionEtape{Comparer les formules}Leurs aires valent $AM\times h:2$ et $MB\times h:2$. Comme $AM=MB$, les deux aires sont égales.\editionRetenir{On n’a pas prouvé que les deux triangles ont la même forme. L’égalité des aires vient des bases et de la hauteur.}
\end{blocEdition}
\editionPage{preuves-4}{Comprendre les preuves}{Démontrer par une formule}{Les arguments ci-dessous éclairent des propriétés utilisées dans le livre.}
\begin{blocEdition}{preuves-4}
\editionEtape{Une affirmation}La somme de trois entiers consécutifs est trois fois celui du milieu.\editionEtape{Nommer}Si le premier est $n$, les suivants sont $n+1$ et $n+2$.\editionEtape{Calculer pour un entier quelconque}$n+(n+1)+(n+2)=3n+3=3(n+1)$. L’affirmation est prouvée pour tous les entiers.\editionEtape{Comprendre la différence avec un test}Vérifier $4+5+6=3\times5$ illustre le résultat mais ne traite qu’un cas. La lettre représente ici tous les cas à la fois.\editionRetenir{Une preuve générale explique pourquoi. Un contre-exemple, au contraire, suffit pour réfuter une affirmation annoncée comme toujours vraie.}
\end{blocEdition}
